\documentclass[12pt,a4paper]{article}

\usepackage[T2A]{fontenc}
\usepackage[utf8]{inputenc}
\usepackage[russian]{babel}
\usepackage{amsmath,amssymb,amsthm,mathtools}
\usepackage{enumitem}
\usepackage{array}
\usepackage{booktabs}
\usepackage{hyperref}

\newtheorem{theorem}{Теорема}
\newtheorem{lemma}{Лемма}
\newtheorem{proposition}{Утверждение}
\newtheorem{definition}{Определение}
\newtheorem{example}{Пример}
\newtheorem{remark}{Замечание}

\DeclareMathOperator{\card}{card}

\title{Билет №14 \\ \small Основы комбинаторного анализа. Метод производящих функций, метод включений и исключений. Генерация и перечисление комбинаторных объектов. Примеры применения. (ИТМО) \\ Основы комбинаторного анализа. Метод производящих функций, метод включений и исключений. Примеры применения. (СпБГУ)}
\author{Сериков Владислав Александрович}
\date{5 мая 2026}


\begin{document}

\maketitle
\tableofcontents
\newpage

\section{Предмет комбинаторного анализа}

Комбинаторный анализ изучает дискретные структуры и методы их подсчёта, построения, кодирования и перечисления. Типичные задачи:
\begin{enumerate}
    \item подсчитать число объектов заданного типа;
    \item построить все объекты заданного типа;
    \item занумеровать объекты, то есть задать биекцию между множеством объектов и множеством номеров;
    \item сгенерировать следующий объект в некотором порядке;
    \item случайно сгенерировать объект с заданным распределением;
    \item доказать тождества между комбинаторными величинами.
\end{enumerate}

\begin{definition}
Комбинаторным классом называется множество дискретных объектов, обычно разбитое на конечные подмножества по размеру:
\[
\mathcal{A} = \bigsqcup_{n \ge 0} \mathcal{A}_n,
\]
где $\mathcal{A}_n$ --- множество объектов размера $n$.
Число
\[
a_n = |\mathcal{A}_n|
\]
называется числом объектов размера $n$.
\end{definition}

Основные комбинаторные методы:
\begin{itemize}
    \item правило суммы;
    \item правило произведения;
    \item биективный метод;
    \item метод рекуррентных соотношений;
    \item метод производящих функций;
    \item метод включений и исключений;
    \item методы генерации и перечисления объектов.
\end{itemize}

\section{Базовые правила подсчёта}

\subsection{Правило суммы}

\begin{theorem}[правило суммы]
Пусть конечное множество $A$ представлено в виде дизъюнктного объединения:
\[
A = A_1 \sqcup A_2 \sqcup \dots \sqcup A_k.
\]
Тогда
\[
|A| = |A_1| + |A_2| + \dots + |A_k|.
\]
\end{theorem}

\begin{proof}
Так как множества $A_i$ попарно не пересекаются, каждый элемент $A$ принадлежит ровно одному из множеств $A_i$. Поэтому при сложении $|A_1|,\dots,|A_k|$ каждый элемент $A$ учитывается ровно один раз. Следовательно,
\[
|A| = \sum_{i=1}^k |A_i|.
\]
\end{proof}

\begin{example}
Сколько существует натуральных чисел от $1$ до $100$, делящихся на $2$ или на $5$, если множества не пересекаются?

Если отдельно считать чётные числа, не делящиеся на $5$, и числа, делящиеся на $5$, то получим дизъюнктное разбиение:
\[
A = A_1 \sqcup A_2.
\]
Далее применяем правило суммы.
\end{example}

\subsection{Правило произведения}

\begin{theorem}[правило произведения]
Пусть объект строится в $k$ последовательных этапов. Если на первом этапе есть $n_1$ вариантов, на втором --- $n_2$ вариантов, и так далее, причём число вариантов на каждом этапе не зависит от предыдущих выборов, то всего объектов:
\[
n_1 n_2 \dots n_k.
\]
\end{theorem}

\begin{proof}
Для $k=2$ каждому объекту соответствует упорядоченная пара $(a,b)$, где $a$ --- выбор на первом этапе, $b$ --- выбор на втором. Для каждого из $n_1$ выборов первого этапа существует $n_2$ выборов второго этапа. Поэтому всего $n_1n_2$ пар.

Для произвольного $k$ доказательство проводится индукцией. Если утверждение верно для $k-1$, то первые $k-1$ этапов можно выполнить
\[
n_1n_2\dots n_{k-1}
\]
способами, а затем для каждого такого выбора есть $n_k$ вариантов последнего этапа. Следовательно, всего
\[
n_1n_2\dots n_{k-1}n_k.
\]
\end{proof}

\begin{example}
Количество двоичных строк длины $n$ равно
\[
2^n,
\]
поскольку на каждой из $n$ позиций можно независимо выбрать $0$ или $1$.
\end{example}

\section{Размещения, перестановки, сочетания}

\subsection{Перестановки}

\begin{definition}
Перестановкой множества из $n$ элементов называется упорядоченное расположение всех его элементов.
\end{definition}

Число перестановок:
\[
P_n = n!.
\]

\begin{proof}
На первое место можно поставить любой из $n$ элементов, на второе --- любой из оставшихся $n-1$, далее $n-2$ и так далее. По правилу произведения:
\[
P_n = n(n-1)(n-2)\dots 1 = n!.
\]
\end{proof}

\subsection{Размещения}

\begin{definition}
Размещением из $n$ элементов по $k$ называется упорядоченный набор из $k$ различных элементов, выбранных из $n$.
\end{definition}

Число размещений:
\[
A_n^k = n(n-1)\dots(n-k+1) = \frac{n!}{(n-k)!}.
\]

\subsection{Сочетания}

\begin{definition}
Сочетанием из $n$ элементов по $k$ называется неупорядоченный набор из $k$ различных элементов, выбранных из $n$.
\end{definition}

Число сочетаний:
\[
C_n^k = \binom{n}{k} = \frac{n!}{k!(n-k)!}.
\]

\begin{proof}
Размещение из $n$ по $k$ можно получить так:
\begin{enumerate}
    \item выбрать $k$ элементов из $n$;
    \item упорядочить выбранные $k$ элементов.
\end{enumerate}
Первый шаг можно выполнить $\binom{n}{k}$ способами, второй --- $k!$ способами. Поэтому
\[
A_n^k = \binom{n}{k}k!.
\]
Но
\[
A_n^k = \frac{n!}{(n-k)!}.
\]
Следовательно,
\[
\binom{n}{k} = \frac{n!}{k!(n-k)!}.
\]
\end{proof}

\section{Биномиальные коэффициенты}

\subsection{Биномиальная теорема}

\begin{theorem}[биномиальная теорема]
Для любого целого $n \ge 0$ имеет место тождество
\[
(x+y)^n = \sum_{k=0}^{n} \binom{n}{k} x^k y^{n-k}.
\]
\end{theorem}

\begin{proof}
Рассмотрим произведение:
\[
(x+y)^n = \underbrace{(x+y)(x+y)\dots(x+y)}_{n \text{ раз}}.
\]
При раскрытии скобок каждый одночлен получается выбором из каждой скобки либо $x$, либо $y$.

Чтобы получить одночлен $x^k y^{n-k}$, нужно выбрать $x$ ровно из $k$ скобок, а $y$ --- из остальных $n-k$ скобок. Число таких выборов равно $\binom{n}{k}$.

Следовательно, коэффициент при $x^k y^{n-k}$ равен $\binom{n}{k}$, и
\[
(x+y)^n = \sum_{k=0}^n \binom{n}{k}x^k y^{n-k}.
\]
\end{proof}

\subsection{Основные тождества}

\[
\binom{n}{k} = \binom{n}{n-k}.
\]

\[
\binom{n}{k} = \binom{n-1}{k} + \binom{n-1}{k-1}.
\]

\begin{proof}
Докажем тождество Паскаля:
\[
\binom{n}{k} = \binom{n-1}{k} + \binom{n-1}{k-1}.
\]
Пусть есть множество из $n$ элементов и выделен один элемент $a$.

Количество $k$-элементных подмножеств равно $\binom{n}{k}$. Разобьём эти подмножества на два класса:
\begin{enumerate}
    \item подмножества, не содержащие $a$;
    \item подмножества, содержащие $a$.
\end{enumerate}
Если подмножество не содержит $a$, то его нужно выбрать из оставшихся $n-1$ элементов:
\[
\binom{n-1}{k}.
\]
Если подмножество содержит $a$, то нужно выбрать ещё $k-1$ элемент из оставшихся $n-1$:
\[
\binom{n-1}{k-1}.
\]
По правилу суммы:
\[
\binom{n}{k} = \binom{n-1}{k} + \binom{n-1}{k-1}.
\]
\end{proof}

\section{Мультимножества и сочетания с повторениями}

\begin{definition}
Сочетанием с повторениями из $n$ типов элементов по $k$ называется выбор $k$ элементов из $n$ типов, где элементы одного типа могут повторяться.
\end{definition}

Число сочетаний с повторениями:
\[
\binom{n+k-1}{k}.
\]

\begin{theorem}
Количество неотрицательных целочисленных решений уравнения
\[
x_1+x_2+\dots+x_n = k
\]
равно
\[
\binom{n+k-1}{k}.
\]
\end{theorem}

\begin{proof}
Каждому решению уравнения соответствует строка из $k$ звёздочек и $n-1$ перегородок:
\[
\underbrace{* \dots *}_{x_1}
|
\underbrace{* \dots *}_{x_2}
|
\dots
|
\underbrace{* \dots *}_{x_n}.
\]
Всего в строке $k+n-1$ позиций: $k$ звёздочек и $n-1$ перегородок. Чтобы задать решение, достаточно выбрать позиции для $k$ звёздочек:
\[
\binom{k+n-1}{k}.
\]
И наоборот, любая такая строка однозначно определяет решение. Следовательно, число решений равно
\[
\binom{n+k-1}{k}.
\]
\end{proof}

\begin{example}
Сколько решений имеет уравнение
\[
x_1+x_2+x_3 = 5,\qquad x_i \ge 0?
\]
Ответ:
\[
\binom{3+5-1}{5} = \binom{7}{5} = 21.
\]
\end{example}

\section{Производящие функции}

\subsection{Обыкновенные производящие функции}

\begin{definition}
Обыкновенной производящей функцией последовательности $(a_n)_{n\ge 0}$ называется формальный степенной ряд
\[
A(x) = \sum_{n\ge 0} a_n x^n.
\]
\end{definition}

Здесь $x$ рассматривается как формальная переменная. Поэтому вопросы сходимости обычно не важны: производящая функция является алгебраическим способом хранения последовательности.

\begin{example}
Последовательности $a_n=1$ соответствует производящая функция
\[
A(x)=1+x+x^2+\dots=\frac{1}{1-x}.
\]
\end{example}

\begin{example}
Последовательности $a_n=n$ соответствует производящая функция
\[
\sum_{n\ge 0} n x^n = \frac{x}{(1-x)^2}.
\]
Действительно,
\[
\frac{1}{1-x}=\sum_{n\ge 0}x^n.
\]
Дифференцируя формально:
\[
\frac{1}{(1-x)^2}=\sum_{n\ge 1} n x^{n-1}.
\]
Умножая на $x$, получаем:
\[
\frac{x}{(1-x)^2}=\sum_{n\ge 1} n x^n.
\]
\end{example}

\subsection{Основные операции с производящими функциями}

Пусть
\[
A(x)=\sum_{n\ge 0}a_nx^n,\qquad B(x)=\sum_{n\ge 0}b_nx^n.
\]

\subsubsection{Сложение}

Если
\[
C(x)=A(x)+B(x),
\]
то
\[
C(x)=\sum_{n\ge 0}(a_n+b_n)x^n.
\]

\subsubsection{Умножение}

Если
\[
C(x)=A(x)B(x),
\]
то коэффициент при $x^n$ равен
\[
c_n = \sum_{k=0}^n a_k b_{n-k}.
\]
Это называется свёрткой последовательностей.

\begin{theorem}[правило произведения для производящих функций]
Пусть объект класса $\mathcal{C}$ получается как упорядоченная пара объектов:
\[
C=(A,B),
\]
где $A \in \mathcal{A}$, $B \in \mathcal{B}$, а размер аддитивен:
\[
|C|=|A|+|B|.
\]
Если
\[
A(x)=\sum_{n\ge 0}a_nx^n,\qquad B(x)=\sum_{n\ge 0}b_nx^n,
\]
то производящая функция класса $\mathcal{C}$ равна
\[
C(x)=A(x)B(x).
\]
\end{theorem}

\begin{proof}
Число объектов класса $\mathcal{C}$ размера $n$ равно числу способов выбрать объект $A$ размера $k$ и объект $B$ размера $n-k$ для некоторого $k$:
\[
c_n=\sum_{k=0}^n a_kb_{n-k}.
\]
С другой стороны, коэффициент при $x^n$ в произведении
\[
A(x)B(x)
\]
равен именно
\[
\sum_{k=0}^n a_kb_{n-k}.
\]
Следовательно,
\[
C(x)=A(x)B(x).
\]
\end{proof}

\subsection{Извлечение коэффициентов}

Обозначение
\[
[x^n]A(x)
\]
означает коэффициент при $x^n$ в ряде $A(x)$.

Например:
\[
[x^n]\frac{1}{1-x}=1,
\]
\[
[x^n]\frac{1}{(1-x)^m}=\binom{n+m-1}{m-1}.
\]

\subsection{Производящие функции и задачи о размене}

\begin{example}
Сколько способов набрать сумму $10$ монетами достоинства $1$, $2$ и $5$?

Для монет достоинства $1$ возможные вклады в сумму:
\[
0,1,2,3,\dots
\]
Производящая функция:
\[
1+x+x^2+x^3+\dots=\frac{1}{1-x}.
\]

Для монет достоинства $2$:
\[
1+x^2+x^4+x^6+\dots=\frac{1}{1-x^2}.
\]

Для монет достоинства $5$:
\[
1+x^5+x^{10}+x^{15}+\dots=\frac{1}{1-x^5}.
\]

Общая производящая функция:
\[
F(x)=\frac{1}{(1-x)(1-x^2)(1-x^5)}.
\]

Искомое число равно:
\[
[x^{10}]F(x).
\]

Перечислим вручную по числу монет достоинства $5$.

Если $c=0$, нужно решить
\[
a+2b=10.
\]
Тогда $b=0,1,2,3,4,5$, всего $6$ решений.

Если $c=1$, нужно решить
\[
a+2b=5.
\]
Тогда $b=0,1,2$, всего $3$ решения.

Если $c=2$, нужно решить
\[
a+2b=0.
\]
Одно решение.

Итого:
\[
6+3+1=10.
\]
\end{example}

\subsection{Производящие функции для ограниченных выборов}

Если объект типа $i$ можно взять $0,1,2,\dots,m_i$ раз, то соответствующий множитель:
\[
1+x+x^2+\dots+x^{m_i}
=
\frac{1-x^{m_i+1}}{1-x}.
\]

Если объект типа $i$ можно взять любое число раз, то множитель:
\[
1+x+x^2+\dots=\frac{1}{1-x}.
\]

Если объект типа $i$ можно взять только $0$ или $1$ раз, то множитель:
\[
1+x.
\]

\begin{example}
Сколько решений имеет уравнение
\[
x_1+x_2+x_3=7,
\]
если
\[
0\le x_1\le 2,\qquad 0\le x_2\le 3,\qquad x_3\ge 0?
\]

Производящая функция:
\[
F(x)=(1+x+x^2)(1+x+x^2+x^3)\frac{1}{1-x}.
\]

Искомое число:
\[
[x^7]F(x).
\]

Так как $x_1+x_2\le 5$, для каждого выбора $x_1,x_2$ однозначно находится
\[
x_3=7-x_1-x_2.
\]
Значит число решений равно числу пар:
\[
x_1\in\{0,1,2\},\qquad x_2\in\{0,1,2,3\}.
\]
Всего:
\[
3\cdot 4=12.
\]
\end{example}

\section{Экспоненциальные производящие функции}

\begin{definition}
Экспоненциальной производящей функцией последовательности $(a_n)_{n\ge 0}$ называется ряд
\[
A(x)=\sum_{n\ge 0}a_n\frac{x^n}{n!}.
\]
\end{definition}

Экспоненциальные производящие функции удобны для подсчёта размеченных объектов, то есть объектов, построенных на множестве меток $\{1,\dots,n\}$.

\begin{example}
Число перестановок $n$ элементов равно $n!$. Экспоненциальная производящая функция:
\[
\sum_{n\ge 0}n!\frac{x^n}{n!}
=
\sum_{n\ge 0}x^n
=
\frac{1}{1-x}.
\]
\end{example}

\begin{example}
Число подмножеств множества из $n$ элементов равно $2^n$. Экспоненциальная производящая функция:
\[
\sum_{n\ge 0}2^n\frac{x^n}{n!}=e^{2x}.
\]
\end{example}

\section{Решение рекуррентных соотношений производящими функциями}

\subsection{Числа Фибоначчи}

Рассмотрим последовательность:
\[
F_0=0,\qquad F_1=1,\qquad F_n=F_{n-1}+F_{n-2},\quad n\ge 2.
\]

Пусть
\[
F(x)=\sum_{n\ge 0}F_nx^n.
\]

Тогда
\[
F(x)=F_0+F_1x+\sum_{n\ge 2}F_nx^n.
\]
Подставим рекуррентное соотношение:
\[
F(x)=x+\sum_{n\ge 2}(F_{n-1}+F_{n-2})x^n.
\]
Разделим сумму:
\[
F(x)=x+\sum_{n\ge 2}F_{n-1}x^n+\sum_{n\ge 2}F_{n-2}x^n.
\]
Преобразуем:
\[
\sum_{n\ge 2}F_{n-1}x^n
=
x\sum_{n\ge 2}F_{n-1}x^{n-1}
=
x\sum_{m\ge 1}F_mx^m
=
xF(x),
\]
так как $F_0=0$.

Аналогично:
\[
\sum_{n\ge 2}F_{n-2}x^n
=
x^2\sum_{m\ge 0}F_mx^m
=
x^2F(x).
\]

Следовательно,
\[
F(x)=x+xF(x)+x^2F(x).
\]
Отсюда
\[
F(x)(1-x-x^2)=x,
\]
\[
F(x)=\frac{x}{1-x-x^2}.
\]

Таким образом, обыкновенная производящая функция чисел Фибоначчи:
\[
\boxed{
F(x)=\frac{x}{1-x-x^2}.
}
\]

\section{Метод включений и исключений}

\subsection{Формула для объединения множеств}

\begin{theorem}[принцип включений и исключений]
Пусть $A_1,A_2,\dots,A_n$ --- конечные множества. Тогда
\[
\left|\bigcup_{i=1}^n A_i\right|
=
\sum_{i=1}^n |A_i|
-
\sum_{1\le i<j\le n}|A_i\cap A_j|
+
\sum_{1\le i<j<k\le n}|A_i\cap A_j\cap A_k|
-\dots
+
(-1)^{n+1}|A_1\cap\dots\cap A_n|.
\]
Иначе:
\[
\left|\bigcup_{i=1}^n A_i\right|
=
\sum_{\varnothing\neq I\subseteq \{1,\dots,n\}}
(-1)^{|I|+1}
\left|\bigcap_{i\in I}A_i\right|.
\]
\end{theorem}

\begin{proof}
Рассмотрим произвольный элемент $x$. Пусть он принадлежит ровно $r$ из множеств $A_1,\dots,A_n$.

Если $r=0$, то $x$ не принадлежит объединению, и в правой части он не учитывается.

Пусть $r\ge 1$. Тогда вклад элемента $x$ в правую часть равен
\[
\binom{r}{1}-\binom{r}{2}+\binom{r}{3}-\dots+(-1)^{r+1}\binom{r}{r}.
\]
Нужно показать, что этот вклад равен $1$.

Из биномиальной теоремы:
\[
(1-1)^r = \sum_{k=0}^r \binom{r}{k}(-1)^k=0.
\]
Следовательно,
\[
\sum_{k=1}^r \binom{r}{k}(-1)^k=-1.
\]
Умножая на $-1$, получаем:
\[
\sum_{k=1}^r (-1)^{k+1}\binom{r}{k}=1.
\]
Значит каждый элемент объединения учитывается ровно один раз, а элементы вне объединения не учитываются. Поэтому формула верна.
\end{proof}

\subsection{Форма для дополнения}

Пусть $U$ --- универсальное множество, а $A_i$ --- множества объектов, обладающих нежелательным свойством $i$.

Тогда число объектов, не обладающих ни одним из свойств $A_1,\dots,A_n$, равно
\[
\left|U\setminus\bigcup_{i=1}^n A_i\right|
=
|U|
-
\left|\bigcup_{i=1}^n A_i\right|.
\]

По принципу включений и исключений:
\[
\left|U\setminus\bigcup_{i=1}^n A_i\right|
=
\sum_{I\subseteq \{1,\dots,n\}}
(-1)^{|I|}
\left|\bigcap_{i\in I}A_i\right|,
\]
где при $I=\varnothing$ считается:
\[
\bigcap_{i\in \varnothing}A_i=U.
\]

\subsection{Число сюръекций}

\begin{theorem}
Количество сюръективных отображений из $n$-элементного множества в $k$-элементное множество равно
\[
\sum_{j=0}^{k}(-1)^j\binom{k}{j}(k-j)^n.
\]
\end{theorem}

\begin{proof}
Пусть
\[
U
\]
--- множество всех отображений из $n$-элементного множества $X$ в $k$-элементное множество $Y$.

Тогда
\[
|U|=k^n.
\]

Отображение не является сюръективным тогда и только тогда, когда существует элемент $y\in Y$, который не имеет прообраза.

Для каждого $i\in Y$ обозначим через $A_i$ множество отображений, не принимающих значение $i$.

Тогда число сюръекций равно:
\[
\left|U\setminus\bigcup_{i\in Y}A_i\right|.
\]

Если выбрано множество $I\subseteq Y$ размера $j$, то отображения из
\[
\bigcap_{i\in I}A_i
\]
не принимают ни одного значения из $I$. Поэтому каждое из $n$ значений можно выбрать из оставшихся $k-j$ элементов. Значит
\[
\left|\bigcap_{i\in I}A_i\right|=(k-j)^n.
\]

По формуле включений и исключений:
\[
\left|U\setminus\bigcup_{i\in Y}A_i\right|
=
\sum_{I\subseteq Y}(-1)^{|I|}(k-|I|)^n.
\]
Группируя по $j=|I|$, получаем:
\[
\sum_{j=0}^{k}(-1)^j\binom{k}{j}(k-j)^n.
\]
\end{proof}

\begin{example}
Сколько сюръекций существует из множества $\{1,2,3\}$ в множество $\{a,b\}$?

По формуле:
\[
2^3-\binom{2}{1}1^3+\binom{2}{2}0^3
=
8-2=6.
\]

Действительно, всего отображений $8$. Не сюръективны только два постоянных отображения.
\end{example}

\subsection{Беспорядки}

\begin{definition}
Беспорядком, или дерранжировкой, называется перестановка $\pi$ множества $\{1,\dots,n\}$ без неподвижных точек:
\[
\pi(i)\neq i
\quad \text{для всех } i.
\]
\end{definition}

\begin{theorem}
Число беспорядков $D_n$ равно
\[
D_n
=
n!\sum_{j=0}^{n}\frac{(-1)^j}{j!}.
\]
\end{theorem}

\begin{proof}
Пусть $U=S_n$ --- множество всех перестановок множества $\{1,\dots,n\}$.
Тогда
\[
|U|=n!.
\]

Для каждого $i$ определим:
\[
A_i=\{\pi\in S_n:\pi(i)=i\}.
\]
То есть $A_i$ --- множество перестановок, фиксирующих элемент $i$.

Нужно найти:
\[
\left|U\setminus\bigcup_{i=1}^n A_i\right|.
\]

Если выбрано множество $I\subseteq \{1,\dots,n\}$ размера $j$, то перестановки из
\[
\bigcap_{i\in I}A_i
\]
фиксируют все элементы из $I$. Остальные $n-j$ элементов можно переставить произвольно. Поэтому
\[
\left|\bigcap_{i\in I}A_i\right|=(n-j)!.
\]

По включениям и исключениям:
\[
D_n
=
\sum_{I\subseteq \{1,\dots,n\}}
(-1)^{|I|}(n-|I|)!.
\]
Группируя по $j=|I|$, получаем:
\[
D_n
=
\sum_{j=0}^{n}(-1)^j\binom{n}{j}(n-j)!.
\]
Но
\[
\binom{n}{j}(n-j)!
=
\frac{n!}{j!(n-j)!}(n-j)!
=
\frac{n!}{j!}.
\]
Следовательно,
\[
D_n
=
n!\sum_{j=0}^{n}\frac{(-1)^j}{j!}.
\]
\end{proof}

\begin{example}
Для $n=4$:
\[
D_4=4!\left(1-1+\frac{1}{2!}-\frac{1}{3!}+\frac{1}{4!}\right)
=
24\left(\frac{1}{2}-\frac{1}{6}+\frac{1}{24}\right)
=
24\cdot \frac{9}{24}
=
9.
\]
\end{example}

\section{Связь включений и исключений с производящими функциями}

Метод включений и исключений часто можно интерпретировать через производящие функции. Особенно полезна идея: запретить некоторую конфигурацию можно вычитанием вкладов объектов, содержащих эту конфигурацию.

\begin{example}
Найдём число неотрицательных решений:
\[
x_1+x_2+x_3=10,
\]
при ограничениях
\[
x_1\le 4,\qquad x_2\le 4,\qquad x_3\le 4.
\]

Без ограничений число решений:
\[
\binom{10+3-1}{3-1}=\binom{12}{2}=66.
\]

Обозначим через $A_i$ множество решений, где $x_i\ge 5$.

Тогда нужно найти:
\[
66-|A_1\cup A_2\cup A_3|.
\]

Считаем $|A_1|$. Если $x_1\ge 5$, положим
\[
y_1=x_1-5\ge 0.
\]
Тогда
\[
y_1+x_2+x_3=5,
\]
и число решений:
\[
\binom{5+3-1}{2}=\binom{7}{2}=21.
\]
Таких множеств три, поэтому вклад:
\[
3\cdot 21=63.
\]

Для пересечения $A_i\cap A_j$ две переменные не меньше $5$. Например,
\[
x_1\ge 5,\quad x_2\ge 5.
\]
Положив
\[
y_1=x_1-5,\quad y_2=x_2-5,
\]
получим:
\[
y_1+y_2+x_3=0.
\]
Одно решение. Пар переменных:
\[
\binom{3}{2}=3.
\]
Тройное пересечение невозможно, так как сумма была бы не меньше $15$.

Искомое число:
\[
66-63+3=6.
\]

С помощью производящих функций та же задача записывается так:
\[
[x^{10}](1+x+x^2+x^3+x^4)^3.
\]
\end{example}

\section{Генерация и перечисление комбинаторных объектов}

\subsection{Понятия генерации, ранжирования и анранжирования}

\begin{definition}
Генерацией комбинаторных объектов называется алгоритмическое построение всех объектов заданного класса.
\end{definition}

\begin{definition}
Перечислением называется задание порядка на объектах и последовательная выдача объектов в этом порядке.
\end{definition}

\begin{definition}
Ранжированием называется вычисление номера объекта в фиксированном порядке.
\end{definition}

\begin{definition}
Анранжированием называется восстановление объекта по его номеру.
\end{definition}

Часто используются следующие порядки:
\begin{itemize}
    \item лексикографический порядок;
    \item ко-лексикографический порядок;
    \item порядок Грея, когда соседние объекты отличаются минимально;
    \item порядок по битовым маскам.
\end{itemize}

\section{Генерация подмножеств}

\subsection{Генерация с помощью битовых масок}

Пусть дано множество
\[
\{1,2,\dots,n\}.
\]
Каждому подмножеству соответствует двоичная маска длины $n$:
\[
b_1b_2\dots b_n,
\]
где
\[
b_i =
\begin{cases}
1, & \text{если } i \text{ входит в подмножество},\\
0, & \text{иначе}.
\end{cases}
\]

Всего масок:
\[
2^n.
\]

\subsubsection*{Алгоритм генерации всех подмножеств}

\begin{enumerate}
    \item Перебрать числа
    \[
    m=0,1,\dots,2^n-1.
    \]
    \item Для каждого числа $m$ рассмотреть его двоичную запись длины $n$.
    \item Если $i$-й бит равен $1$, включить элемент $i$ в текущее подмножество.
    \item Вывести полученное подмножество.
\end{enumerate}

\subsubsection*{Минимальный пример}

Пусть
\[
n=3.
\]
Тогда маски и подмножества:

\[
\begin{array}{c|c}
\text{Маска} & \text{Подмножество}\\
\hline
000 & \varnothing\\
001 & \{3\}\\
010 & \{2\}\\
011 & \{2,3\}\\
100 & \{1\}\\
101 & \{1,3\}\\
110 & \{1,2\}\\
111 & \{1,2,3\}
\end{array}
\]

Если нумеровать биты справа налево, то маска $101$ соответствует подмножеству $\{1,3\}$.

\subsection{Генерация в порядке Грея}

Код Грея --- последовательность битовых строк, в которой соседние строки отличаются ровно в одном бите.

\subsubsection*{Рекурсивное построение}

Обозначим через $G_n$ список кодов Грея длины $n$.

\[
G_1=(0,1).
\]

Для построения $G_n$:
\begin{enumerate}
    \item взять список $G_{n-1}$ и приписать к каждому слову слева $0$;
    \item взять список $G_{n-1}$ в обратном порядке и приписать к каждому слову слева $1$;
    \item объединить два полученных списка.
\end{enumerate}

То есть:
\[
G_n = 0G_{n-1},\,1\operatorname{rev}(G_{n-1}).
\]

\subsubsection*{Пример}

\[
G_1:
\quad
0,\ 1.
\]

\[
G_2:
\quad
00,\ 01,\ 11,\ 10.
\]

\[
G_3:
\quad
000,\ 001,\ 011,\ 010,\ 110,\ 111,\ 101,\ 100.
\]

Соседние подмножества отличаются добавлением или удалением одного элемента.

\section{Генерация сочетаний}

\subsection{Лексикографический порядок сочетаний}

Сочетание из $n$ по $k$ можно представить возрастающей последовательностью:
\[
1\le a_1<a_2<\dots<a_k\le n.
\]

\subsubsection*{Алгоритм получения следующего сочетания}

Пусть текущее сочетание:
\[
(a_1,a_2,\dots,a_k).
\]

\begin{enumerate}
    \item Найти максимальный индекс $i$, такой что
    \[
    a_i < n-k+i.
    \]
    \item Если такого $i$ нет, текущее сочетание последнее.
    \item Увеличить $a_i$ на $1$:
    \[
    a_i := a_i+1.
    \]
    \item Для всех $j>i$ положить:
    \[
    a_j := a_{j-1}+1.
    \]
    \item Полученная последовательность --- следующее сочетание.
\end{enumerate}

\subsubsection*{Почему условие имеет вид $a_i<n-k+i$}

В последнем сочетании:
\[
(n-k+1,n-k+2,\dots,n).
\]
Поэтому максимальное возможное значение $a_i$ равно:
\[
n-k+i.
\]

\subsubsection*{Минимальный пример}

Пусть
\[
n=5,\qquad k=3.
\]
Сочетания в лексикографическом порядке:

\[
\begin{array}{c}
(1,2,3)\\
(1,2,4)\\
(1,2,5)\\
(1,3,4)\\
(1,3,5)\\
(1,4,5)\\
(2,3,4)\\
(2,3,5)\\
(2,4,5)\\
(3,4,5)
\end{array}
\]

Пусть текущее сочетание:
\[
(1,3,5).
\]
Ищем справа налево индекс $i$:
\[
a_3=5,\quad n-k+3=5,
\]
увеличить нельзя.

\[
a_2=3,\quad n-k+2=4,
\]
увеличить можно. Значит $i=2$.

Делаем:
\[
a_2:=4,
\]
а затем:
\[
a_3:=a_2+1=5.
\]
Получаем:
\[
(1,4,5).
\]

\section{Генерация перестановок}

\subsection{Лексикографический следующий элемент}

Пусть дана перестановка
\[
(a_1,a_2,\dots,a_n).
\]

\subsubsection*{Алгоритм следующей перестановки}

\begin{enumerate}
    \item Найти максимальный индекс $i$, такой что
    \[
    a_i<a_{i+1}.
    \]
    \item Если такого $i$ нет, перестановка является последней.
    \item Найти максимальный индекс $j>i$, такой что
    \[
    a_j>a_i.
    \]
    \item Поменять местами $a_i$ и $a_j$.
    \item Развернуть хвост:
    \[
    a_{i+1},a_{i+2},\dots,a_n.
    \]
\end{enumerate}

\subsubsection*{Идея алгоритма}

Хвост после позиции $i$ является невозрастающим. Чтобы получить минимально большую перестановку:
\begin{enumerate}
    \item нужно немного увеличить префикс;
    \item затем сделать хвост минимально возможным, то есть отсортировать его по возрастанию.
\end{enumerate}
Так как хвост был невозрастающим, после обмена достаточно его развернуть.

\subsubsection*{Минимальный пример}

Пусть текущая перестановка:
\[
(1,3,5,4,2).
\]

Шаг 1. Ищем справа налево индекс $i$, где есть рост:
\[
3<5.
\]
Значит
\[
i=2,\quad a_i=3.
\]

Шаг 2. Ищем справа налево элемент больше $3$:
\[
4>3.
\]
Значит
\[
j=4,\quad a_j=4.
\]

Шаг 3. Меняем:
\[
(1,4,5,3,2).
\]

Шаг 4. Разворачиваем хвост после позиции $2$:
\[
(5,3,2)\mapsto(2,3,5).
\]

Получаем:
\[
(1,4,2,3,5).
\]

\section{Ранжирование и анранжирование перестановок}

\subsection{Факториальная система счисления}

Любую перестановку можно закодировать последовательностью чисел:
\[
c_1,c_2,\dots,c_n,
\]
где $c_i$ --- число ещё не использованных элементов, меньших $a_i$.

Тогда ранг перестановки в лексикографическом порядке:
\[
\operatorname{rank}(a_1,\dots,a_n)
=
\sum_{i=1}^{n} c_i (n-i)!.
\]

Если ранги начинаются с $0$, то первая перестановка имеет ранг $0$.

\subsubsection*{Пример}

Рассмотрим перестановку:
\[
(2,4,1,3).
\]

Список доступных элементов сначала:
\[
[1,2,3,4].
\]

\[
a_1=2.
\]
Меньше $2$ среди доступных только $1$, значит
\[
c_1=1.
\]

Удаляем $2$:
\[
[1,3,4].
\]

\[
a_2=4.
\]
Меньше $4$ среди доступных $1,3$, значит
\[
c_2=2.
\]

Удаляем $4$:
\[
[1,3].
\]

\[
a_3=1.
\]
Меньше $1$ среди доступных нет, значит
\[
c_3=0.
\]

\[
a_4=3.
\]
\[
c_4=0.
\]

Ранг:
\[
1\cdot 3!+2\cdot 2!+0\cdot 1!+0\cdot 0!
=
6+4=10.
\]

\section{Примеры применения методов}

\subsection{Подсчёт слов с ограничениями}

\begin{example}
Сколько двоичных строк длины $n$ не содержат двух единиц подряд?

Обозначим это число через $a_n$.

Рассмотрим последнюю позицию строки:
\begin{enumerate}
    \item если строка оканчивается на $0$, то первые $n-1$ символов образуют любую допустимую строку длины $n-1$;
    \item если строка оканчивается на $10$, точнее последняя единица стоит в конце, то перед ней обязательно стоит $0$ или это начало строки. Удобнее рассматривать строки, оканчивающиеся на $1$: тогда предыдущий символ обязан быть $0$, а первые $n-2$ символов образуют допустимую строку.
\end{enumerate}

Получаем рекурсию:
\[
a_n=a_{n-1}+a_{n-2}.
\]

Начальные значения:
\[
a_0=1,\qquad a_1=2.
\]

Следовательно,
\[
a_n=F_{n+2},
\]
где $F_m$ --- числа Фибоначчи с $F_0=0,F_1=1$.

Производящая функция:
\[
A(x)=\sum_{n\ge 0}a_nx^n.
\]
Из рекурсии:
\[
A(x)=\frac{1+x}{1-x-x^2}.
\]
\end{example}

\subsection{Подсчёт целочисленных решений с верхними ограничениями}

\begin{example}
Найти число решений:
\[
x_1+x_2+x_3+x_4=12,
\]
где
\[
0\le x_i\le 5.
\]

Без верхних ограничений:
\[
\binom{12+4-1}{3}=\binom{15}{3}=455.
\]

Пусть $A_i$ --- множество решений с $x_i\ge 6$.

Если одно ограничение нарушено:
\[
x_i=y_i+6,
\]
получаем:
\[
y_i+\sum_{j\ne i}x_j=6.
\]
Число решений:
\[
\binom{6+4-1}{3}=\binom{9}{3}=84.
\]
Таких $i$ четыре, вклад:
\[
4\cdot 84.
\]

Если два ограничения нарушены:
\[
x_i=y_i+6,\qquad x_j=y_j+6.
\]
Тогда остаётся сумма $0$, значит одно решение. Пар:
\[
\binom{4}{2}=6.
\]

Три ограничения нарушены быть не могут, так как сумма была бы не меньше $18$.

По включениям и исключениям:
\[
455-4\cdot84+6=455-336+6=125.
\]

Через производящие функции:
\[
[x^{12}](1+x+x^2+x^3+x^4+x^5)^4=125.
\]
\end{example}

\subsection{Разбиения числа}

\begin{definition}
Разбиением числа $n$ называется представление $n$ в виде суммы положительных целых слагаемых без учёта порядка.
\end{definition}

Например, разбиения числа $4$:
\[
4,\quad 3+1,\quad 2+2,\quad 2+1+1,\quad 1+1+1+1.
\]
Значит
\[
p(4)=5.
\]

Производящая функция числа разбиений:
\[
\sum_{n\ge 0}p(n)x^n
=
\prod_{m=1}^{\infty}\frac{1}{1-x^m}.
\]

\begin{proof}
Слагаемое $m$ может входить в разбиение $0,1,2,\dots$ раз. Поэтому его вклад в производящую функцию:
\[
1+x^m+x^{2m}+x^{3m}+\dots=\frac{1}{1-x^m}.
\]
Так как выбор кратностей для разных $m$ независим, общая производящая функция равна произведению:
\[
\prod_{m=1}^{\infty}\frac{1}{1-x^m}.
\]
Коэффициент при $x^n$ равен числу способов выбрать кратности слагаемых так, чтобы общая сумма была $n$, то есть числу разбиений $p(n)$.
\end{proof}

\section{Числа Каталана}

\subsection{Определение и первые значения}

\begin{definition}
Числами Каталана называется последовательность $(C_n)_{n\ge 0}$, задаваемая формулой
\[
C_n=\frac{1}{n+1}\binom{2n}{n}.
\]
\end{definition}

Первые значения:
\[
C_0=1,\quad
C_1=1,\quad
C_2=2,\quad
C_3=5,\quad
C_4=14,\quad C_5=42,\quad C_6=132.
\]

Числа Каталана возникают во многих задачах комбинаторики. Они считают, например:
\begin{itemize}
    \item правильные скобочные последовательности длины $2n$;
    \item пути Дика длины $2n$;
    \item способы триангуляции выпуклого $(n+2)$-угольника;
    \item количество бинарных деревьев с $n$ внутренними вершинами;
    \item количество способов корректно расставить скобки в произведении из $n+1$ множителей.
\end{itemize}

\subsection{Правильные скобочные последовательности}

\begin{definition}
Правильной скобочной последовательностью длины $2n$ называется строка из $n$ открывающих и $n$ закрывающих скобок такая, что в каждом её префиксе число открывающих скобок не меньше числа закрывающих.
\end{definition}

Например, для $n=3$ существуют $5$ правильных скобочных последовательностей:
\[
((())),\quad (()()),\quad (())(),\quad ()(()),\quad ()()().
\]

\begin{theorem}
Количество правильных скобочных последовательностей длины $2n$ равно числу Каталана:
\[
C_n=\frac{1}{n+1}\binom{2n}{n}.
\]
\end{theorem}

\begin{proof}
Рассмотрим все последовательности из $n$ открывающих и $n$ закрывающих скобок. Их количество равно
\[
\binom{2n}{n},
\]
поскольку нужно выбрать позиции для $n$ открывающих скобок.

Назовём последовательность плохой, если в некотором её префиксе закрывающих скобок становится больше, чем открывающих.

Посчитаем число плохих последовательностей. Для этого построим биекцию между плохими последовательностями с $n$ открывающими и $n$ закрывающими скобками и произвольными последовательностями из $n-1$ открывающих и $n+1$ закрывающих скобок.

Пусть дана плохая последовательность. Рассмотрим первый момент, когда число закрывающих скобок становится на $1$ больше числа открывающих. То есть рассмотрим минимальный префикс, в котором
\[
\#)=\#(+1.
\]
В этом префиксе заменим все открывающие скобки на закрывающие, а все закрывающие --- на открывающие. Остальную часть последовательности оставим без изменений.

Пусть в выбранном префиксе было $a$ открывающих и $a+1$ закрывающих скобок. После замены в нём станет $a+1$ открывающих и $a$ закрывающих.

Изначально во всей последовательности было $n$ открывающих и $n$ закрывающих скобок. После такой операции число открывающих увеличится на $1$, а число закрывающих уменьшится на $1$ относительно префикса. Однако если использовать противоположную конвенцию отражения, плохие последовательности с $n$ открывающими и $n$ закрывающими скобками биективно соответствуют последовательностям с $n+1$ закрывающей и $n-1$ открывающей скобкой.

Удобно записать это в терминах шагов пути. Пусть открывающая скобка соответствует шагу вверх $+1$, а закрывающая --- шагу вниз $-1$. Правильная последовательность соответствует пути из $(0,0)$ в $(2n,0)$, который никогда не опускается ниже оси.

Плохой путь впервые достигает уровня $-1$. Отразим часть пути до первого достижения уровня $-1$ относительно горизонтальной оси. После отражения путь будет иметь на один шаг вверх меньше и на один шаг вниз больше, то есть будет состоять из
\[
n-1
\]
шагов вверх и
\[
n+1
\]
шагов вниз.

Обратно, любой путь из $n-1$ шагов вверх и $n+1$ шагов вниз обязательно заканчивается на уровне $-2$, поэтому обязательно впервые достигает уровня $-1$. Отразив начальный участок до первого достижения уровня $-1$, получим плохой путь с $n$ шагами вверх и $n$ шагами вниз.

Следовательно, число плохих последовательностей равно
\[
\binom{2n}{n-1},
\]
так как нужно выбрать позиции для $n-1$ открывающих скобок среди $2n$ позиций.

Поэтому число правильных последовательностей равно
\[
\binom{2n}{n}-\binom{2n}{n-1}.
\]
Преобразуем:
\[
\binom{2n}{n-1}
=
\binom{2n}{n}\frac{n}{n+1}.
\]
Действительно,
\[
\frac{\binom{2n}{n-1}}{\binom{2n}{n}}
=
\frac{\frac{(2n)!}{(n-1)!(n+1)!}}{\frac{(2n)!}{n!n!}}
=
\frac{n!n!}{(n-1)!(n+1)!}
=
\frac{n}{n+1}.
\]
Следовательно,
\[
\binom{2n}{n}-\binom{2n}{n-1}
=
\binom{2n}{n}\left(1-\frac{n}{n+1}\right)
=
\frac{1}{n+1}\binom{2n}{n}.
\]
Значит число правильных скобочных последовательностей равно
\[
C_n=\frac{1}{n+1}\binom{2n}{n}.
\]
\end{proof}

\subsection{Рекуррентное соотношение для чисел Каталана}

\begin{theorem}
Числа Каталана удовлетворяют рекуррентному соотношению
\[
C_0=1,
\]
\[
C_{n+1}=\sum_{k=0}^{n}C_kC_{n-k},\qquad n\ge 0.
\]
\end{theorem}

\begin{proof}
Докажем рекурсию на модели правильных скобочных последовательностей.

Рассмотрим правильную скобочную последовательность длины $2(n+1)$. Она начинается с открывающей скобки. Пусть закрывающая скобка, парная к первой открывающей, стоит так, что внутри первой пары находится правильная скобочная последовательность длины $2k$, а после этой пары --- правильная скобочная последовательность длины $2(n-k)$.

То есть вся последовательность имеет вид
\[
(S)T,
\]
где $S$ --- правильная скобочная последовательность длины $2k$, а $T$ --- правильная скобочная последовательность длины $2(n-k)$.

Для фиксированного $k$ последовательность $S$ можно выбрать $C_k$ способами, а последовательность $T$ --- $C_{n-k}$ способами. По правилу произведения получаем
\[
C_kC_{n-k}
\]
вариантов.

Так как $k$ может принимать значения от $0$ до $n$, по правилу суммы:
\[
C_{n+1}=\sum_{k=0}^{n}C_kC_{n-k}.
\]
\end{proof}

\begin{example}
Вычислим $C_4$ по рекурсии.

Известно:
\[
C_0=1,\quad C_1=1,\quad C_2=2,\quad C_3=5.
\]
Тогда
\[
C_4=C_0C_3+C_1C_2+C_2C_1+C_3C_0.
\]
Следовательно,
\[
C_4=1\cdot 5+1\cdot 2+2\cdot 1+5\cdot 1=14.
\]
\end{example}

\subsection{Производящая функция чисел Каталана}

Пусть
\[
C(x)=\sum_{n\ge 0}C_nx^n.
\]

Из рекурсии
\[
C_{n+1}=\sum_{k=0}^{n}C_kC_{n-k}
\]
получаем уравнение для производящей функции:
\[
C(x)=1+xC(x)^2.
\]

\begin{theorem}
Обыкновенная производящая функция чисел Каталана равна
\[
C(x)=\frac{1-\sqrt{1-4x}}{2x}.
\]
\end{theorem}

\begin{proof}
Из рекуррентного соотношения:
\[
C_0=1,\qquad C_{n+1}=\sum_{k=0}^{n}C_kC_{n-k}
\]
получаем:
\[
C(x)-1
=
\sum_{n\ge 0}C_{n+1}x^{n+1}.
\]
Подставим рекурсию:
\[
C(x)-1
=
\sum_{n\ge 0}\left(\sum_{k=0}^{n}C_kC_{n-k}\right)x^{n+1}.
\]
Вынесем $x$:
\[
C(x)-1
=
x\sum_{n\ge 0}\left(\sum_{k=0}^{n}C_kC_{n-k}\right)x^{n}.
\]
Внутренняя сумма является коэффициентом при $x^n$ в произведении $C(x)^2$. Поэтому
\[
C(x)-1=xC(x)^2.
\]
Следовательно,
\[
C(x)=1+xC(x)^2.
\]

Это квадратное уравнение относительно $C(x)$:
\[
xC(x)^2-C(x)+1=0.
\]
Решая его, получаем:
\[
C(x)=\frac{1\pm\sqrt{1-4x}}{2x}.
\]

Нужно выбрать знак. Так как
\[
C(0)=C_0=1,
\]
производящая функция должна быть формальным степенным рядом с конечным свободным членом.

Рассмотрим две ветви.

Для знака плюс:
\[
\frac{1+\sqrt{1-4x}}{2x}
\]
имеет особенность порядка $x^{-1}$ при $x=0$, поскольку числитель стремится к $2$.

Для знака минус:
\[
\frac{1-\sqrt{1-4x}}{2x}
\]
числитель обращается в ноль при $x=0$, и выражение имеет конечный предел. Поэтому
\[
C(x)=\frac{1-\sqrt{1-4x}}{2x}.
\]
\end{proof}

\subsection{Вывод явной формулы через производящую функцию}

\begin{theorem}
Для всех $n\ge 0$
\[
C_n=\frac{1}{n+1}\binom{2n}{n}.
\]
\end{theorem}

\begin{proof}
Из уравнения
\[
C(x)=1+xC(x)^2
\]
можно получить явную формулу через разложение
\[
C(x)=\frac{1-\sqrt{1-4x}}{2x}.
\]

Используем обобщённую биномиальную формулу:
\[
(1+z)^\alpha
=
\sum_{m\ge 0}\binom{\alpha}{m}z^m.
\]
При $\alpha=\frac12$ и $z=-4x$:
\[
\sqrt{1-4x}
=
\sum_{m\ge 0}\binom{1/2}{m}(-4x)^m.
\]
Тогда
\[
1-\sqrt{1-4x}
=
-\sum_{m\ge 1}\binom{1/2}{m}(-4x)^m.
\]
Следовательно,
\[
C(x)=
-\frac{1}{2x}\sum_{m\ge 1}\binom{1/2}{m}(-4x)^m.
\]
Положим $m=n+1$. Тогда коэффициент при $x^n$ равен
\[
C_n
=
-\frac12 \binom{1/2}{n+1}(-4)^{n+1}.
\]

Преобразование этого выражения даёт
\[
C_n=\frac{1}{n+1}\binom{2n}{n}.
\]

Можно также получить формулу из доказанного ранее подсчёта правильных скобочных последовательностей:
\[
C_n=
\binom{2n}{n}-\binom{2n}{n-1}
=
\frac{1}{n+1}\binom{2n}{n}.
\]
\end{proof}

\subsection{Триангуляции выпуклого многоугольника}

\begin{definition}
Триангуляцией выпуклого многоугольника называется разбиение его на треугольники непересекающимися диагоналями.
\end{definition}

\begin{theorem}
Количество триангуляций выпуклого $(n+2)$-угольника равно числу Каталана
\[
C_n.
\]
\end{theorem}

\begin{proof}
Обозначим через $T_n$ число триангуляций выпуклого $(n+2)$-угольника. Для треугольника имеем
\[
T_1=1.
\]
Удобнее положить
\[
T_0=1,
\]
что соответствует пустой или вырожденной конфигурации.

Зафиксируем одну сторону многоугольника, например сторону между вершинами $0$ и $n+1$. В любой триангуляции эта сторона принадлежит ровно одному треугольнику. Пусть третья вершина этого треугольника имеет номер $k+1$, где
\[
0\le k\le n-1.
\]

Тогда выбранный треугольник разбивает многоугольник на две меньшие части:
\begin{itemize}
    \item левую часть, которая триангулируется $T_k$ способами;
    \item правую часть, которая триангулируется $T_{n-1-k}$ способами.
\end{itemize}

Для фиксированного $k$ число триангуляций равно
\[
T_kT_{n-1-k}.
\]
Суммируя по $k$, получаем:
\[
T_n=\sum_{k=0}^{n-1}T_kT_{n-1-k}.
\]
При $T_0=1$ это та же рекурсия, что и для чисел Каталана:
\[
C_n=\sum_{k=0}^{n-1}C_kC_{n-1-k}.
\]
Следовательно,
\[
T_n=C_n.
\]
Значит количество триангуляций выпуклого $(n+2)$-угольника равно $C_n$.
\end{proof}

\begin{example}
Для пятиугольника имеем $n+2=5$, значит $n=3$. Количество триангуляций:
\[
C_3=5.
\]
\end{example}

\subsection{Бинарные деревья}

\begin{definition}
Полным бинарным деревом называется корневое дерево, в котором каждая внутренняя вершина имеет ровно двух потомков.
\end{definition}

\begin{theorem}
Количество полных бинарных деревьев с $n$ внутренними вершинами равно
\[
C_n.
\]
\end{theorem}

\begin{proof}
Обозначим через $B_n$ число полных бинарных деревьев с $n$ внутренними вершинами.

Если $n=0$, имеется одно дерево, состоящее из одного листа:
\[
B_0=1.
\]

Если $n\ge 1$, то корень является внутренней вершиной. Его левое поддерево содержит $k$ внутренних вершин, а правое --- $n-1-k$ внутренних вершин для некоторого $k$ от $0$ до $n-1$.

Левое поддерево можно выбрать $B_k$ способами, правое --- $B_{n-1-k}$ способами. Поэтому
\[
B_n=\sum_{k=0}^{n-1}B_kB_{n-1-k}.
\]

Это рекурсия Каталана с тем же начальным условием:
\[
B_0=1.
\]
Следовательно,
\[
B_n=C_n.
\]
\end{proof}

\subsection{Пути Дика}

\begin{definition}
Путём Дика полудлины $n$ называется путь на квадратной решётке из точки $(0,0)$ в точку $(2n,0)$, состоящий из шагов
\[
U=(1,1),\qquad D=(1,-1),
\]
который никогда не опускается ниже оси $Ox$.
\end{definition}

Правильные скобочные последовательности длины $2n$ находятся в биекции с путями Дика:
\[
( \leftrightarrow U,
\qquad
) \leftrightarrow D.
\]
Поэтому число путей Дика полудлины $n$ равно
\[
C_n.
\]

\subsection{Алгоритм генерации правильных скобочных последовательностей}

\subsubsection*{Идея}

Нужно построить все строки длины $2n$, содержащие $n$ открывающих и $n$ закрывающих скобок, такие что ни в одном префиксе число закрывающих скобок не превосходит число открывающих.

\subsubsection*{Алгоритм с возвратом}

Пусть:
\begin{itemize}
    \item $open$ --- число уже поставленных открывающих скобок;
    \item $close$ --- число уже поставленных закрывающих скобок;
    \item $s$ --- текущий префикс.
\end{itemize}

Рекурсивная процедура:

\begin{enumerate}
    \item Если
    \[
    |s|=2n,
    \]
    вывести $s$ и завершить текущую ветвь.
    \item Если
    \[
    open<n,
    \]
    можно добавить открывающую скобку:
    \[
    s:=s+\texttt{(}.
    \]
    Затем рекурсивно продолжить.
    \item Если
    \[
    close<open,
    \]
    можно добавить закрывающую скобку:
    \[
    s:=s+\texttt{)}.
    \]
    Затем рекурсивно продолжить.
\end{enumerate}

\subsubsection*{Минимальный пример для $n=3$}

Начинаем с пустой строки:
\[
s=\varepsilon,\quad open=0,\quad close=0.
\]

Первый символ может быть только открывающей скобкой:
\[
(
\]

Дальше возможны ветви:
\[
(( \quad \text{или} \quad ()
\]

В результате алгоритм выдаст:
\[
((())),\quad (()()),\quad (())(),\quad ()(()),\quad ()()().
\]

\subsubsection*{Корректность алгоритма}

Алгоритм никогда не строит неправильный префикс, поскольку закрывающая скобка добавляется только при условии
\[
close<open.
\]
Это означает, что после добавления закрывающей скобки число закрывающих скобок не превысит число открывающих.

Алгоритм никогда не использует больше $n$ открывающих скобок, поскольку открывающая скобка добавляется только при условии
\[
open<n.
\]
Закрывающих скобок также не может стать больше $n$, поскольку
\[
close\le open\le n.
\]

Каждая правильная скобочная последовательность может быть построена этим алгоритмом единственным образом: на каждом шаге алгоритм выбирает ровно тот символ, который стоит на соответствующей позиции данной последовательности. Следовательно, алгоритм порождает все и только правильные скобочные последовательности.

\section{Применения комбинаторных методов в математике и науке}

\subsection{Комбинаторика как универсальный язык дискретных структур}

Комбинаторные методы применяются во всех областях, где требуется:
\begin{itemize}
    \item считать конечные или дискретные объекты;
    \item оценивать число возможных состояний системы;
    \item анализировать алгоритмы;
    \item строить модели случайных дискретных процессов;
    \item изучать структуры графов, деревьев, слов, разбиений, перестановок;
    \item доказывать тождества и получать рекуррентные формулы.
\end{itemize}

Производящие функции, принцип включений и исключений, рекурсии, биномиальные коэффициенты и методы генерации объектов являются базовыми инструментами не только в дискретной математике, но и в теории вероятностей, алгебре, анализе алгоритмов, статистической физике, биоинформатике и теории информации.

\subsection{Применения в алгебре}

\subsubsection{Биномиальные коэффициенты и многочлены}

Биномиальные коэффициенты естественно возникают при раскрытии степеней:
\[
(x+y)^n=\sum_{k=0}^{n}\binom{n}{k}x^ky^{n-k}.
\]
Они также появляются в формулах для конечных разностей, интерполяции и разложениях многочленов.

\begin{example}
Оператор конечной разности определяется как
\[
\Delta f(x)=f(x+1)-f(x).
\]
Для многочленов степени $n$ повторное применение конечной разности приводит к суммам с биномиальными коэффициентами:
\[
\Delta^n f(x)
=
\sum_{k=0}^{n}(-1)^{n-k}\binom{n}{k}f(x+k).
\]
\end{example}

\subsubsection{Симметрические многочлены}

Комбинаторные объекты используются для описания симметрических многочленов. Например, элементарный симметрический многочлен степени $k$ от переменных $x_1,\dots,x_n$ задаётся формулой
\[
e_k(x_1,\dots,x_n)
=
\sum_{1\le i_1<i_2<\dots<i_k\le n}
x_{i_1}x_{i_2}\dots x_{i_k}.
\]
Число слагаемых в нём равно
\[
\binom{n}{k}.
\]

\subsection{Применения в теории графов}

Теория графов является одной из центральных областей применения комбинаторики. Графы моделируют сети, отношения, маршруты, зависимости и конфликты.

\subsubsection{Подсчёт графов}

На множестве из $n$ вершин существует
\[
\binom{n}{2}
\]
возможных неориентированных рёбер без петель. Каждое ребро либо присутствует, либо отсутствует. Поэтому число простых неориентированных графов на фиксированном множестве из $n$ вершин равно
\[
2^{\binom{n}{2}}.
\]

Для ориентированных графов без петель между каждой упорядоченной парой различных вершин может быть дуга. Таких пар:
\[
n(n-1).
\]
Следовательно, число ориентированных графов без петель равно
\[
2^{n(n-1)}.
\]

\subsubsection{Деревья}

\begin{theorem}[формула Кэли]
Количество помеченных деревьев на множестве вершин
\[
\{1,2,\dots,n\}
\]
равно
\[
n^{n-2}.
\]
\end{theorem}

\begin{proof}
Доказательство проводится с помощью кодов Прюфера.

Каждому помеченному дереву на $n$ вершинах сопоставим последовательность длины $n-2$ над алфавитом $\{1,\dots,n\}$.

Алгоритм построения кода Прюфера:
\begin{enumerate}
    \item Пока в дереве больше двух вершин, выбрать лист с наименьшей меткой.
    \item Записать метку его единственного соседа.
    \item Удалить выбранный лист из дерева.
    \item Повторять процесс до тех пор, пока не останутся две вершины.
\end{enumerate}

На каждом из $n-2$ шагов записывается одна метка, поэтому получается последовательность длины $n-2$.

Обратно, по любой последовательности
\[
(p_1,p_2,\dots,p_{n-2})
\]
можно восстановить дерево:
\begin{enumerate}
    \item Рассмотреть множество текущих вершин $\{1,\dots,n\}$.
    \item На шаге $i$ выбрать наименьшую метку $v$, которая не встречается в оставшейся части кода
    \[
    p_i,p_{i+1},\dots,p_{n-2}.
    \]
    \item Добавить ребро
    \[
    \{v,p_i\}.
    \]
    \item Удалить $v$ из множества текущих вершин.
    \item После обработки всех элементов кода останутся две вершины; соединить их ребром.
\end{enumerate}

Эти две процедуры взаимно обратны. Следовательно, существует биекция между помеченными деревьями на $n$ вершинах и последовательностями длины $n-2$ над $n$-элементным алфавитом.

Число таких последовательностей равно
\[
n^{n-2}.
\]
Значит число помеченных деревьев равно
\[
n^{n-2}.
\]
\end{proof}

\begin{example}
Для $n=4$ число помеченных деревьев:
\[
4^{4-2}=16.
\]
\end{example}

\subsection{Применения в теории вероятностей}

Комбинаторика лежит в основе классической вероятности на конечных пространствах элементарных исходов.

Если все исходы равновозможны, то вероятность события $A$ равна
\[
P(A)=\frac{|A|}{|\Omega|}.
\]

\subsubsection{Гипергеометрическое распределение}

\begin{example}
В урне находится $N$ шаров, из них $M$ белых и $N-M$ чёрных. Из урны без возвращения выбирают $n$ шаров. Найдём вероятность того, что среди выбранных ровно $k$ белых.

Общее число способов выбрать $n$ шаров:
\[
\binom{N}{n}.
\]

Благоприятный выбор:
\begin{itemize}
    \item выбрать $k$ белых из $M$:
    \[
    \binom{M}{k};
    \]
    \item выбрать $n-k$ чёрных из $N-M$:
    \[
    \binom{N-M}{n-k}.
    \]
\end{itemize}

По правилу произведения число благоприятных исходов:
\[
\binom{M}{k}\binom{N-M}{n-k}.
\]

Следовательно,
\[
P(X=k)=
\frac{\binom{M}{k}\binom{N-M}{n-k}}{\binom{N}{n}}.
\]
\end{example}

\subsubsection{Метод включений и исключений в вероятности}

Если $A_1,\dots,A_n$ --- события, то
\[
P\left(\bigcup_{i=1}^{n}A_i\right)
=
\sum_i P(A_i)
-
\sum_{i<j}P(A_i\cap A_j)
+\dots+
(-1)^{n+1}P(A_1\cap\dots\cap A_n).
\]

\begin{example}
Пусть случайно выбирается перестановка $n$ писем по $n$ конвертам. Вероятность того, что ни одно письмо не попадёт в свой конверт, равна
\[
\frac{D_n}{n!}
=
\sum_{j=0}^{n}\frac{(-1)^j}{j!}.
\]
При больших $n$ эта вероятность стремится к
\[
e^{-1}.
\]
\end{example}

\subsection{Применения в анализе алгоритмов}

Комбинаторика используется для оценки:
\begin{itemize}
    \item числа входных данных;
    \item числа возможных состояний;
    \item сложности переборных алгоритмов;
    \item среднего времени работы;
    \item числа сравнений, перестановок, инверсий;
    \item структуры рекурсивных вызовов.
\end{itemize}

\subsubsection{Перебор подмножеств}

Если алгоритм перебирает все подмножества $n$-элементного множества, то он имеет не менее
\[
2^n
\]
итераций. Поэтому такие алгоритмы обычно имеют экспоненциальную сложность.

\begin{example}
Задача о рюкзаке в простом переборном варианте: для каждого из $n$ предметов решить, брать его или нет. Всего вариантов:
\[
2^n.
\]
\end{example}

\subsubsection{Перебор перестановок}

Если алгоритм перебирает все порядки $n$ объектов, то число вариантов равно
\[
n!.
\]

\begin{example}
Наивное решение задачи коммивояжёра перебирает все маршруты. Если фиксировать начальный город, остаётся
\[
(n-1)!
\]
перестановок остальных городов.
\end{example}

\subsubsection{Рекурсии и производящие функции}

Производящие функции помогают получать асимптотику и явные формулы для последовательностей, возникающих в анализе рекурсивных алгоритмов.

Например, если структура рекурсивного алгоритма описывается бинарным деревом, то число возможных форм таких деревьев часто выражается через числа Каталана:
\[
C_n\sim \frac{4^n}{n^{3/2}\sqrt{\pi}}.
\]

\subsection{Применения в информатике и программировании}

\subsubsection{Генерация тестов}

Алгоритмы генерации комбинаторных объектов используются для:
\begin{itemize}
    \item построения всех тестов малой размерности;
    \item проверки корректности программ полным перебором;
    \item стресс-тестирования;
    \item генерации структурных объектов: графов, деревьев, строк, разбиений.
\end{itemize}

\subsubsection{Кодирование и сжатие}

Перечисление и ранжирование позволяют кодировать объект его номером. Если объектов $N$, то для хранения номера нужно примерно
\[
\lceil \log_2 N\rceil
\]
бит.

\begin{example}
Перестановку из $n$ элементов можно закодировать её рангом от $0$ до $n!-1$. Для хранения такого ранга требуется
\[
\lceil \log_2(n!)\rceil
\]
бит.
\end{example}

\subsubsection{Криптография}

Комбинаторика применяется для оценки размера пространства ключей и вероятности успешной атаки полным перебором.

Если ключ имеет длину $m$ бит, то число ключей равно
\[
2^m.
\]
При равномерном выборе ключа вероятность угадать его с одной попытки:
\[
2^{-m}.
\]

\subsection{Применения в теории информации}

Комбинаторные оценки используются в задачах кодирования, передачи и хранения данных.

\subsubsection{Коды исправления ошибок}

Пусть рассматриваются двоичные слова длины $n$. Их всего:
\[
2^n.
\]
Число слов на расстоянии Хэмминга не более $t$ от фиксированного слова равно
\[
\sum_{i=0}^{t}\binom{n}{i}.
\]

Эта величина называется объёмом шара Хэмминга радиуса $t$.

\begin{example}
Если код исправляет до $t$ ошибок, то шары Хэмминга радиуса $t$ вокруг кодовых слов не должны пересекаться. Поэтому если код содержит $M$ слов, то необходимо:
\[
M\sum_{i=0}^{t}\binom{n}{i}\le 2^n.
\]
Это называется границей Хэмминга.
\end{example}

\subsection{Применения в статистической физике}

В статистической физике комбинаторика используется для подсчёта числа микросостояний системы.

\begin{example}
Пусть имеется $N$ частиц, каждая из которых может находиться в одном из двух состояний. Тогда число микросостояний:
\[
2^N.
\]

Если ровно $k$ частиц находятся в первом состоянии, то число микросостояний:
\[
\binom{N}{k}.
\]
\end{example}

Энтропия в статистической физике связана с числом микросостояний $\Omega$ формулой
\[
S=k_B\ln \Omega,
\]
где $k_B$ --- постоянная Больцмана.

\subsection{Применения в биоинформатике}

Комбинаторика применяется при анализе последовательностей ДНК, РНК и белков.

\subsubsection{Число последовательностей}

ДНК-последовательность длины $n$ над алфавитом
\[
\{A,C,G,T\}
\]
может быть выбрана
\[
4^n
\]
способами.

Белковая последовательность длины $n$ над алфавитом из $20$ аминокислот может быть выбрана
\[
20^n
\]
способами.

\subsubsection{Подсчёт слов с ограничениями}

Если требуется посчитать число последовательностей, не содержащих запрещённых подслов, то используются:
\begin{itemize}
    \item рекуррентные соотношения;
    \item конечные автоматы;
    \item матричные производящие функции;
    \item метод включений и исключений.
\end{itemize}

\begin{example}
Количество двоичных строк длины $n$, не содержащих двух единиц подряд, равно
\[
F_{n+2},
\]
где $F_m$ --- числа Фибоначчи.
Такие модели возникают при подсчёте последовательностей с локальными ограничениями.
\end{example}

\subsection{Применения в химии}

В химии комбинаторные методы используются для подсчёта изомеров и различных молекулярных структур.

Например:
\begin{itemize}
    \item графы моделируют молекулы;
    \item вершины соответствуют атомам;
    \item рёбра соответствуют химическим связям;
    \item симметрии молекулы учитываются с помощью групповых действий;
    \item для подсчёта неэквивалентных раскрасок и структур используется лемма Бернсайда и теория Пойа.
\end{itemize}

\begin{example}
Если молекулярный каркас имеет симметрии, то простое число способов заместить атомы нужно делить не напрямую на размер группы симметрий, а учитывать число конфигураций, фиксируемых каждой симметрией. Именно это делает лемма Бернсайда:
\[
\frac{1}{|G|}\sum_{g\in G}|\operatorname{Fix}(g)|.
\]
\end{example}

\subsection{Применения в экономике и социальных науках}

Комбинаторика используется для анализа вариантов выбора, голосований, коалиций и распределений ресурсов.

\subsubsection{Теория голосования}

Если есть $m$ кандидатов, то число возможных строгих ранжирований кандидатов одним избирателем равно
\[
m!.
\]
Если есть $N$ избирателей, то число профилей предпочтений равно
\[
(m!)^N.
\]

\subsubsection{Коалиции}

Для множества из $n$ участников число возможных коалиций равно
\[
2^n.
\]
Если пустая коалиция не учитывается, то:
\[
2^n-1.
\]

\subsection{Применения в машинном обучении и анализе данных}

Комбинаторные методы используются для:
\begin{itemize}
    \item подсчёта числа возможных признаковых подмножеств;
    \item анализа сложности моделей;
    \item построения деревьев решений;
    \item оценки числа разбиений данных;
    \item перебора гиперпараметров;
    \item изучения комбинаторной размерности классов функций.
\end{itemize}

\begin{example}
Если имеется $d$ признаков, то число всех подмножеств признаков равно
\[
2^d.
\]
Поэтому полный перебор всех наборов признаков быстро становится невозможным при росте $d$.
\end{example}

\subsection{Комбинаторные тождества как инструмент доказательства}

Комбинаторные методы часто позволяют доказывать алгебраические тождества через подсчёт одного и того же множества двумя способами.

\begin{example}
Докажем тождество Вандермонда:
\[
\sum_{k=0}^{r}\binom{m}{k}\binom{n}{r-k}
=
\binom{m+n}{r}.
\]

Рассмотрим множество из $m+n$ элементов, разбитое на две группы:
\[
A,\quad |A|=m,
\]
\[
B,\quad |B|=n.
\]
Нужно выбрать $r$ элементов из объединения $A\cup B$.

С одной стороны, это можно сделать
\[
\binom{m+n}{r}
\]
способами.

С другой стороны, можно выбрать $k$ элементов из $A$ и $r-k$ элементов из $B$. Для фиксированного $k$ число способов:
\[
\binom{m}{k}\binom{n}{r-k}.
\]
Суммируя по всем $k$, получаем:
\[
\sum_{k=0}^{r}\binom{m}{k}\binom{n}{r-k}.
\]
Оба выражения считают одно и то же множество, значит
\[
\sum_{k=0}^{r}\binom{m}{k}\binom{n}{r-k}
=
\binom{m+n}{r}.
\]
\end{example}

\subsection{Итоговая схема применения комбинаторики}

При решении прикладной или теоретической задачи полезно пройти следующие шаги:
\begin{enumerate}
    \item определить, какие объекты являются элементарными;
    \item понять, различаются ли объекты порядком, метками, симметриями;
    \item выбрать подходящий метод:
    \begin{itemize}
        \item правило суммы или произведения;
        \item биномиальные коэффициенты;
        \item рекурсии;
        \item производящие функции;
        \item включения и исключения;
        \item лемма Бернсайда;
        \item алгоритмическая генерация;
    \end{itemize}
    \item получить формулу, рекурсию или алгоритм;
    \item проверить результат на малых значениях;
    \item при необходимости оценить асимптотику.
\end{enumerate}

\section{Типовые схемы применения}

\subsection{Когда использовать производящие функции}

Метод производящих функций полезен, если:
\begin{itemize}
    \item нужно посчитать число решений уравнений в неотрицательных целых числах;
    \item есть ограничения на количество объектов каждого типа;
    \item есть рекуррентное соотношение;
    \item объект можно разложить на независимые компоненты;
    \item требуется доказать тождество для последовательностей;
    \item нужно получить явную формулу или асимптотику.
\end{itemize}

Общая схема:
\begin{enumerate}
    \item определить последовательность $a_n$;
    \item записать производящую функцию
    \[
    A(x)=\sum_{n\ge 0}a_nx^n;
    \]
    \item выразить $A(x)$ через элементарные функции или рациональные выражения;
    \item найти нужный коэффициент
    \[
    [x^n]A(x).
    \]
\end{enumerate}

\subsection{Когда использовать включения и исключения}

Метод включений и исключений полезен, если:
\begin{itemize}
    \item нужно посчитать объекты без запрещённых свойств;
    \item проще считать объекты, нарушающие ограничения;
    \item есть пересекающиеся условия;
    \item требуется посчитать сюръекции, беспорядки, раскраски с ограничениями;
    \item возникают условия вида ``не имеет'', ``не содержит'', ``не делится'', ``не совпадает''.
\end{itemize}

Общая схема:
\begin{enumerate}
    \item задать универсальное множество $U$;
    \item определить плохие множества $A_1,\dots,A_n$;
    \item выразить ответ как
    \[
    \left|U\setminus\bigcup_i A_i\right|;
    \]
    \item вычислить размеры пересечений
    \[
    \left|\bigcap_{i\in I}A_i\right|;
    \]
    \item применить формулу:
    \[
    \sum_{I\subseteq\{1,\dots,n\}}(-1)^{|I|}
    \left|\bigcap_{i\in I}A_i\right|.
    \]
\end{enumerate}

\section{Краткая таблица основных формул}

\[
\begin{array}{c|c}
\text{Объект} & \text{Число объектов}\\
\hline
\text{Перестановки } n \text{ элементов} & n!\\
\text{Размещения из } n \text{ по } k & \dfrac{n!}{(n-k)!}\\
\text{Сочетания из } n \text{ по } k & \binom{n}{k}\\
\text{Сочетания с повторениями} & \binom{n+k-1}{k}\\
\text{Подмножества } n\text{-элементного множества} & 2^n\\
\text{Сюръекции } [n]\to[k] &
\sum_{j=0}^k(-1)^j\binom{k}{j}(k-j)^n\\
\text{Беспорядки} &
n!\sum_{j=0}^n\dfrac{(-1)^j}{j!}
\end{array}
\]

\section{Заключение}

Комбинаторный анализ опирается на несколько универсальных идей. Правила суммы и произведения дают базовый аппарат подсчёта. Биномиальные коэффициенты описывают выбор подмножеств и возникают в большом числе задач. Производящие функции позволяют переводить комбинаторные задачи в алгебраические операции со степенными рядами. Метод включений и исключений позволяет считать объекты с запретами и пересекающимися условиями. Алгоритмы генерации, ранжирования и анранжирования позволяют не только подсчитывать объекты, но и эффективно строить их в заданном порядке.

\end{document}
